% !TEX root = Tesi_Corsi_Leonardo_656276.tex
% Il codice sorgente, letto direttamente dai file del progetto. Interlinea
% singola e caratteri non spaziati: con le impostazioni dei capitoli le righe
% di commento andrebbero a capo.
\appendix
\lstset{basicstyle=\footnotesize\ttfamily\linespread{1}\selectfont,
  columns=fullflexible,keepspaces=true,upquote=true}

\chapter{Codice del canale}
\lstinputlisting[language=C,caption={\texttt{src/channel.h}}]{../src/channel.h}

\chapter{Modifiche all'interprete}
Dell'interprete si riportano le parti modificate rispetto a quello di
partenza~\cite{TesiStella}: l'intestazione, con la cella estesa dai canali e
dai due contatori di misura, il controllo degli accessi in memoria con le istruzioni
di \emph{load} e \emph{store} che lo usano, e l'esecuzione delle istruzioni di
comunicazione, di \texttt{ECALL} e delle codifiche non implementate. Prelievo,
decodifica ed esecuzione delle altre istruzioni RV32IM sono quelli
dell'appendice A.2 di~\cite{TesiStella}, a meno della traccia, resa
disattivabile.

\lstinputlisting[language=C,caption={\texttt{src/risc.h}}]{../src/risc.h}
\lstinputlisting[language=C,firstline=153,lastline=224,firstnumber=153,
  caption={\texttt{src/risc.c}, righe 153--224}]{../src/risc.c}
\lstinputlisting[language=C,firstline=439,lastline=500,firstnumber=439,
  caption={\texttt{src/risc.c}, righe 439--500}]{../src/risc.c}

\chapter{Codice della griglia}
\lstinputlisting[language=C,caption={\texttt{src/grid.h}}]{../src/grid.h}
\lstinputlisting[language=C,caption={\texttt{src/grid.c}}]{../src/grid.c}

\chapter{Modello RTL}
\lstinputlisting[language={[System]Verilog},
  caption={\texttt{hw/neso\_defs.svh}}]{../hw/neso_defs.svh}
\lstinputlisting[language={[System]Verilog},
  caption={\texttt{hw/neso\_channel.sv}}]{../hw/neso_channel.sv}
\lstinputlisting[language={[System]Verilog},
  caption={\texttt{hw/neso\_if.sv}}]{../hw/neso_if.sv}
